\Needspace{10\baselineskip}
\subsection{A.158. A compact Lipschitz resolution over every o-minimal real closed field}

\begin{saresult}{A.158}{Proved under stated hypotheses}{A uniform spatial family has a finite compact Lipschitz parameter resolution}
Let \(M\) be an arbitrary o-minimal expansion of a real closed field, let
\(P\subset[0,1]^k\) be a nonempty compact definable set, and let
\[
 F:P\times[0,1]^n\longrightarrow[0,1]^r
\]
be continuous and definable. Use maximum norms. Suppose
\[
 \|F(p,x)-F(p,x')\|_\infty\le K_0\|x-x'\|_\infty
 \quad(p\in P,\ x,x'\in[0,1]^n)
 \tag{158.1}
\]
for some \(K_0\in M_{\ge0}\). There are standard integers
\(J,D,C_\phi,C\), a compact set
\[
 K=\{1,\ldots,J\}\times[0,1]^D,
\]
and a continuous definable surjection \(\phi:K\to P\), such that
\(\phi\) is \(C_\phi\)-Lipschitz and
\[
 (z,x)\longmapsto F(\phi(z),x)
 \quad\hbox{is }C(1+K_0)\hbox{-Lipschitz}.
 \tag{158.2}
\]
The standard integers may depend on the given definable family; no
uniformity across all formulas or all parameter values of an additional
unfixed family is claimed.

Consequently every Hausdorff-continuous definable family of nonempty
compact subsets \(C_p\subset[0,1]^n\), with compact definable parameter
set \(P\subset[0,1]^k\), has a compact Lipschitz resolution with standard
integer projection and Hausdorff constants. Combined with A.157, this
gives relative Hausdorff-Lipschitz regularization on a parameter cube
over every o-minimal real closed field: if the boundary family is
Hausdorff-Lipschitz, there is a definable Lipschitz \(g:[0,1]^k\to[0,1]^k\),
equal to the identity on the boundary, such that \(C_{g(p)}\) is
Hausdorff-Lipschitz. The convex homotopy from the identity to \(g\) fixes
the boundary.
\end{saresult}

\begin{proof}
The only external differentiable input is the generic
\(\mathcal C^1\) boundary-extension claim in the proof of Lemma 2.4 of
Kocel-Cynk--Pawlucki--Valette,
\emph{\(\mathcal C^p\)-parametrization in o-minimal structures},
Canadian Mathematical Bulletin 62 (2019), 99--108,
printed pp.~102--103, DOI \(10.4153/CMB-2018-030-x\).
Its introduction, printed p.~99, explicitly permits any o-minimal
expansion of any real closed field. The claim says that a definable
\(\mathcal C^1\) scalar function on \(W\times(0,1)\), where \(W\) is
open, with bounded derivative in the last variable, extends
\(\mathcal C^1\) to \(W\times[0,1)\) after deleting a closed nowhere
dense subset of the boundary parameter set \(W\). We use this claim,
not an unproved all-field extension of the 2014 Hausdorff-resolution
theorem. We supply the needed version with a nonsmooth exceptional
set explicitly.

\textbf{Generic bounded parameter gradients.}
Let \(U\subset M^d\) be open and definable, and let
\(H:U\times[0,1]^n\to[0,1]^r\) be continuous and definable with the
uniform spatial bound (158.1). Choose an open dense definable subset
\(\mathcal D\subset U\times(0,1)^n\) on which all components of \(H\)
are \(\mathcal C^1\). By dimension of fibres, after removing a closed
nowhere dense subset of \(U\), each fibre
\(\mathcal D_t\) is dense in \((0,1)^n\).
We claim that, after deleting a further closed nowhere dense subset,
all the sets
\[
 \left\{\frac{\partial H_j}{\partial t_i}(t,x):
                   x\in\mathcal D_t\right\}
 \quad(1\le i\le d,\ 1\le j\le r)
 \tag{158.3}
\]
are bounded.

It suffices to rule out an open parameter set \(W\) where one of
these sets is unbounded. Definable choice, followed by a finite
\(\mathcal C^1\) cell decomposition with the parameter coordinates
first, gives, after shrinking \(W\), a definable \(\mathcal C^1\) map
\(\alpha(t,s)\) on \(0<s<\epsilon(t)\), with
\((t,\alpha(t,s))\in\mathcal D\), whose selected parameter derivative
has absolute value tending to infinity as \(s\to0\).
For example, choose its absolute value greater than \(1/s\) before
the cell refinement. The bounded curve \(\alpha(t,s)\) has a
definable limit \(a(t)\in[0,1]^n\); shrink \(W\) so \(a\) is
\(\mathcal C^1\).

After one more finite refinement, one spatial coordinate
\(\alpha_l\) has nonzero \(s\)-derivative of fixed sign and dominates
all the other absolute \(s\)-derivatives. Such a nonzero coordinate
exists, since an eventually constant curve could not carry a
derivative tending to infinity. Write
\[
 v=\varepsilon_l(\alpha_l(t,s)-a_l(t))>0,
 \qquad \varepsilon_l\in\{1,-1\}.
\]
Inverting this coordinate gives a \(\mathcal C^1\) map
\(\widetilde\alpha(t,v)\) for \(0<v<\beta(t)\), with \(\beta>0\)
definable and \(\mathcal C^1\), and
\[
 \left|\frac{\partial\widetilde\alpha_i}{\partial v}\right|
 \le1\quad(1\le i\le n).
\]
Set \(b(t)=\beta(t)/(2(1+\beta(t)))\) and
\(\gamma(t,u)=\widetilde\alpha(t,b(t)u)\), \(0<u<1\).
Then \(|\partial\gamma_i/\partial u|\le1\), the image still lies in
\(\mathcal D\), and the selected parameter derivative still tends
to infinity at \(u=0\). Apply the cited generic boundary-extension
claim to every component of \(\gamma\). They extend
\(\mathcal C^1\) at generic boundary parameters. The function
\(h(t,u)=H_j(t,\gamma(t,u))\) is \(\mathcal C^1\), and
\[
 \left|\frac{\partial h}{\partial u}\right|\le nK_0.
\]
Apply the same claim to \(h\). At a parameter where all these
extensions are \(\mathcal C^1\), the chain rule gives
\[
 \frac{\partial H_j}{\partial t_i}(t,\gamma(t,u))
 =
 \frac{\partial h}{\partial t_i}(t,u)
 -\sum_{\nu=1}^n
   \frac{\partial H_j}{\partial x_\nu}(t,\gamma(t,u))
   \frac{\partial\gamma_\nu}{\partial t_i}(t,u).
 \tag{158.4}
\]
Every term on the right is bounded as \(u\to0\): the spatial
partials have absolute value at most \(K_0\), and the other
partials extend continuously. This contradicts the choice of
\(\gamma\), proving the claim. All bad-parameter sets are definable;
a definable set with empty interior has nowhere dense closure,
which gives the asserted closed exceptional set. When \(n=0\),
there is no spatial supremum, and the assertion follows directly
from generic \(\mathcal C^1\) regularity.

\textbf{A finite norming collection of parameter gradients.}
On a good open parameter set, define
\[
 \mathcal G_t=
 \{e_1,\ldots,e_d\}
 \ \cup
 \{\nabla_tH_j(t,x):x\in\mathcal D_t,\ 1\le j\le r\}.
\]
For each \(t\), this is a bounded definable subset of \(M^d\).
The number
\[
 A(t)=\sup\{|\det(g_1,\ldots,g_d)|:g_i\in\mathcal G_t\}
\]
is finite and at least \(1\). Definable choice supplies an ordered
tuple \(g_1(t),\ldots,g_d(t)\in\mathcal G_t\) with
\[
 |\det(g_1(t),\ldots,g_d(t))|>\tfrac12 A(t).
 \tag{158.5}
\]
For every \(g\in\mathcal G_t\), Cramer's rule then gives
\[
 g=\sum_{i=1}^d c_i g_i(t),\qquad |c_i|\le2.
 \tag{158.6}
\]
Indeed, each numerator determinant replaces one selected vector
by another member of \(\mathcal G_t\), so is at most \(A(t)\).

Choose witnesses for the selected gradients. After a finite
parameter refinement, each selected vector is either a fixed
basis vector \(e_l\), or
\(\nabla_tH_{j_i}(t,\lambda_i(t))\), where
\((t,\lambda_i(t))\in\mathcal D\). Off a lower-dimensional
parameter set the \(\lambda_i\) are \(\mathcal C^1\), and the
finite labels \(j_i,l\) are fixed. In the second case put
\(h_i(t)=H_{j_i}(t,\lambda_i(t))\). In the basis-vector case put
\(\lambda_i(t)=0\) and \(h_i(t)=t_l\).
For a differentiable parameter curve \(t=t(u)\), these choices
give the exact formulas
\[
 \begin{aligned}
 g_i(t)\cdot\partial_{u_\nu}t
 &=\partial_{u_\nu}h_i
   -\nabla_xH_{j_i}(t,\lambda_i(t))
                  \cdot\partial_{u_\nu}\lambda_i
      &&\text{for a function witness},\\
 g_i(t)\cdot\partial_{u_\nu}t
 &=\partial_{u_\nu}t_l
      &&\text{for a basis witness}.
 \end{aligned}
 \tag{158.7}
\]
Thus it suffices to control the finitely many bounded coordinate
functions \(t,\lambda_i,h_i\); no bound on the size of the original
parameter gradients is required.

\textbf{Standard bounds from the finite coordinate clocks.}
We record the quantitative part of A.154 that is used here.
If \(q_0:[0,1]^e\to[0,1]^a\) is continuous and definable, there is
a definable homeomorphism \(\chi\) of the source cube, preserving
every coordinate face, for which \(q=q_0\chi\) is Lipschitz with
a standard integer constant \(L_q\). Apply the coordinate clocks
to \((\operatorname{id},q_0)\). Every component stays in \([0,1]\).
At a stage with \(N\) monotonicity pieces and \(a+e\) components,
the total variation is at most \((a+e)N\), so the new coordinate
bound is at most \(1+(a+e)N\). Previously obtained bounds are
multiplied by at most \(1+4(a+e)N\), by the variation comparison
in A.154. The number of stages and all monotonicity counts are
ordinary finite integers. Starting from the first coordinate
bound therefore gives a standard integer bound. This assertion
does not rely on the field being Archimedean and does not claim
that a field-valued Lipschitz constant can always be rounded to
a standard integer.

\textbf{Resolving one parameter stratum.}
We prove by induction on dimension a slightly stronger statement:
for every definable \(S\subset P\), finitely many compact cube
charts have images in \(\overline S\), cover \(S\), have standard
Lipschitz projection constants, and have pulled \(F\) satisfying
(158.2). Zero-dimensional sets are finite; constant charts work,
with spatial bound \(K_0\). The empty set needs no charts.

Decompose a positive-dimensional \(S\) into finitely many
\(\mathcal C^1\) cells. Lower-dimensional cells are covered by
induction. A cell of dimension \(d\) has a definable
\(\mathcal C^1\) parametrization
\[
 p:(0,1)^d\longrightarrow S_0
\]
which is a diffeomorphism onto that cell. Its image is bounded
in \([0,1]^k\), although its derivatives need not be bounded.
Apply the preceding arguments to \(H(t,x)=F(p(t),x)\).
Remove the lower-dimensional exceptional parameter set and
refine so all selected witnesses and their finite labels have
the stated regularity. The images under \(p\) of the removed
sets have dimension less than \(d\), so induction will cover
them.

On each remaining open parameter piece \(W\), form the bounded
graph
\[
 \Gamma_\Phi=
 \{(t,p(t),\lambda_1(t),\ldots,\lambda_d(t),
                    h_1(t),\ldots,h_d(t)):t\in W\}.
 \tag{158.8}
\]
All its coordinates belong to unit cubes. Triangulate the compact
closure \(\overline{\Gamma_\Phi}\) compatibly with
\(\Gamma_\Phi\). Take the closures of all simplices whose relative
interiors are mapped into \(\Gamma_\Phi\). These finitely many
closed simplices cover \(\overline{\Gamma_\Phi}\): a compatible
finite triangulation represents the graph by a union of open
simplices, and its closure is the union of their closed
simplices. Parametrize each such closed \(e\)-simplex by a fixed
PL homeomorphism from \([0,1]^e\), and compose with the definable
triangulation map. Apply the preceding coordinate-clock
construction to this finite vector map. The resulting chart
\[
 q:[0,1]^e\longrightarrow\overline{\Gamma_\Phi}
\]
is Lipschitz with a standard integer constant \(L_q\), has the
same image as before, and maps its cube interior into
\(\Gamma_\Phi\). Let \(\phi_q\) be its \(p\)-coordinate.
Its image lies in \(\overline{S_0}\subset P\), and it is
\(L_q\)-Lipschitz.

\textbf{Controlling the entire pulled spatial family.}
Set \(\widetilde F(u,x)=F(\phi_q(u),x)\). This is continuous on
the compact product cube. On the cube interior the coordinate
relations (158.8) hold. The chart \(q\) is \(\mathcal C^1\) on
a dense open subset. Moreover, for every good \(t\), the spatial
exceptional fibre \((0,1)^n\setminus\mathcal D_t\) has dimension
less than \(n\). Thus its pullback under \(u\mapsto t(u)\) has
dimension less than \(e+n\), by dimension of fibres. On a dense
open subset of the product cube the chart and the composed
function are therefore \(\mathcal C^1\), and all formulas
(158.6)--(158.7) apply.

Each coordinate derivative of \(q\) is bounded by \(L_q\).
For a selected function witness, (158.7) gives
\[
 |g_i(t(u))\cdot\partial_{u_\nu}t|
       \le L_q(1+nK_0);
\]
for a basis witness the bound \(L_q\) suffices. Hence, for every
scalar component of the pulled family,
\[
 \left|\frac{\partial\widetilde F_j}{\partial u_\nu}\right|
 \le 2d L_q(1+nK_0),\qquad
 \left|\frac{\partial\widetilde F_j}{\partial x_\mu}\right|
 \le K_0
 \tag{158.9}
\]
on that dense open subset. This remains valid for a chart
simplex of dimension \(e<d\); the spatial exceptional-fibre
bound was imposed for every parameter of \(W\), not merely
for generic parameters of the ambient \(d\)-space.

For clarity, bounded derivatives on such a dense open set imply
the required global Lipschitz bound here. A continuous definable
scalar function on a convex cube, with all coordinate derivatives
bounded on a dense open subset, is Lipschitz with the sum of
those bounds. Given an interior segment direction, project the
lower-dimensional bad set along that direction. Outside a
lower-dimensional set of parallel offsets, a parallel segment
meets the bad set in finitely many points, by dimension of
fibres. Split at those finitely many points, apply the Mean
Value Theorem on the good open intervals, and use continuity
at the splitting points. Arbitrarily small such parallel
translations exist. Pass to the given segment by continuity,
and then to boundary endpoints by interior approximation.
Every partition in this argument is finite and definable;
neither integration nor an ordinary finite mesh is used.
Applying this argument componentwise to (158.9) proves a
Lipschitz bound
\[
 C_q(1+K_0)
\]
with a standard integer \(C_q\), for example any standard
integer at least \(2edL_q(n+1)+n+1\).

This produces the required charts over \(W\), and the
lower-dimensional induction covers all removed parameter
sets. All resulting images remain in \(\overline S\).
For \(S=P\), their union is therefore exactly \(P\).
Pad lower-dimensional source cubes by unused coordinates to
one common dimension \(D\), and put the finitely many charts
on the distinct integer tags \(1,\ldots,J\). Between different
tags the source distance is at least \(1\), whereas all
projection and function values lie in unit cubes. The
within-chart constants thus give global standard constants
\(C_\phi,C\) as in (158.2). This proves the first assertion.

\textbf{Compact-set families and the relative consequence.}
For \(C_p\subset[0,1]^n\), define
\[
 F(p,x)=\operatorname{dist}_\infty(x,C_p),
 \qquad x\in[0,1]^n.
\]
The distance is a definable minimum, belongs to \([0,1]\), is
\(1\)-Lipschitz in \(x\), and is jointly continuous by Hausdorff
continuity of the original family. For nonempty compact
\(A,B\subset[0,1]^n\),
\[
 d_H(A,B)=
 \sup_{x\in[0,1]^n}
 |\operatorname{dist}_\infty(x,A)
             -\operatorname{dist}_\infty(x,B)|.
 \tag{158.10}
\]
The upper inequality is the triangle inequality; the reverse
inequality is obtained by evaluating at a point of \(A\) or
\(B\) attaining the directed Hausdorff distance. These
maxima exist by definable compactness. Applying (158.2) with
\(K_0=1\) and then (158.10) gives a compact resolution with
standard Hausdorff and projection constants.

The initial metric hypothesis of A.157 is now fulfilled over
every o-minimal real closed field. Apply its relative
construction, including the affine tagged boundary section
and the normalized finite coordinate clocks. This supplies
the stated \(g\) and its relative convex homotopy. The empty
parameter set has the unique empty resolution; zero-dimensional
parameter cubes and separately allowed empty spatial families
have their immediate trivial versions.
\end{proof}

\paragraph{Scope.}
This proof supplies an all-field compact resolution for continuous
families. It does not assert that the real-field 2014 theorem was
originally stated in this generality, nor does it assert an all-field
classification of arbitrary external Hausdorff limits. Its constants
are standard because the derivative-witness graph is bounded in unit
cubes before the finite variation clocks are applied. No derivative
bound on the original parameter chart, no ordinary Smale theorem,
and no Archimedean subdivision argument is used.
